%%%%%%%%%%%%%%%%%%%%%%%%%%%%%%%%%%%%%%%%%%%%%%%%%%%%%%%%%%%%%%%%%%%%%%
% Lớp 10 | Hàm số bậc hai và đồ thị | Hàm số và đồ thị
% Dạng 0D3?1-7: Toán thực tế  về hàm số
% Số câu: 9  (sinh tự động bởi tach_cau.py; từ nay thầy sửa trực tiếp file này)
%%%%%%%%%%%%%%%%%%%%%%%%%%%%%%%%%%%%%%%%%%%%%%%%%%%%%%%%%%%%%%%%%%%%%%

%%%%%%%%%%%%%%%% TRẮC NGHIỆM 4 PHƯƠNG ÁN %%%%%%%%%%%%%%%%

% nguồn: MaHoaNH_TracNghiem_Lop10
\begin{ex}%[0D3H1-7]%[ma:dcc12548]
	\immini{Hàm số mô tả sự phụ thuộc  của $y$ (\textit{số tiền phải trả, tính bằng nghìn đồng}) vào $x$ (\textit{lượng điện tiêu thụ}, \textit{tính bằng} $kWh$) trên từng khoảng giá trị $x$ được cho bằng công thức như hình bên.\\
	Ông An tháng thứ nhất sài hết 90 kWh điện, tháng thứ hai sài hết 120 kWh điện. Theo công thức  bên thì số tiền ông An phải trả trong 2 tháng gần nhất với kết quả nào sau đây?}{\hspace{1cm}
	$y=\heva{&1{,}678x &\text{ nếu }& 0 \le x \le 50\\&1{,}734x-2{,}8 &\text{ nếu }& 50 < x \le 100\\&2{,}014x-30{,}8 &\text{ nếu }& 100 < x \le 200\\&2{,}536x-135{,}2 &\text{ nếu }& 200 < x \le 300\\&2{,}834 x-224{,}6 &\text{ nếu }& 300 < x \le  400\\&2{,}927 x-261{,}8 &\text{ nếu }& x>400.}$}
	 \choice
		{\True $364$ nghìn  đồng}
		{$361$ nghìn đồng} 
		{$392$ nghìn  đồng}
		{$352$ nghìn  đồng}
	\loigiai{ 
		Số tiền phải trả là $1{,}734\cdot 90-2{,}8  + 2{,}014\cdot -30{,}8=364{,}14$ (nghìn đồng).
	 }
\end{ex}


%%%%%%%%%%%%%%%% ĐÚNG/SAI %%%%%%%%%%%%%%%%

% nguồn: MaHoaNH_DungSai_Lop10
\begin{ex}%[0D3H1-7]%[ma:5ec24b7d]
	Một hiệu chuyên cho thuê xe máy niêm yết giá như sau: Giá thuê xe là $110$ nghìn đồng một ngày cho ba ngày đầu tiên và $80$ nghìn đồng cho mỗi ngày tiếp theo. Mỗi kết quả dưới  đây đúng hay sai?
	 \choiceTF
		 {\True Tổng số tiền phải trả T (nghìn đồng) theo số ngày $x$ được xác định bởi $T(x)=80x+90$ nếu $x>3$}
		{$T(2)=230$} 
		{\True Với số tiền là $2$ triệu đồng thì thuê xe tối đa trong $22$ ngày liên tiếp} 
		{$T(x)$ là hàm đồng  biến $\left(0;+\infty\right)$}
	 \loigiai{
		\begin{itemchoice}
			\itemch Ta có $T(x)=\heva{& 110x &  \text{nếu}\,0\le x\le 3\\& 80x+90 & \text{nếu}\, x>3.}$
			\itemch  $T(2)=110\cdot 2=220$ nghìn đồng.
			\itemch Với $0\le x \le 3$  thì $ 110x\le 330$.\\
			Với $2$ triệu đồng số tiền khách còn sau khi thuê $3$ ngày đầu  là $2000-330=1670$ nghìn đồng.\\
			Xét $80x+90\le 1670\Leftrightarrow  x\le 19{,}75$.\\
			Với số tiền là $2$ triệu đồng thì thuê xe  tối đa trong $3+19=22$ ngày liên tiếp
			\itemch $T(x)$ là hàm số đồng  biến.
		\end{itemchoice}
	} 
\end{ex}

% nguồn: MaHoaNH_DungSai_Lop10
\begin{ex}%[0D3H1-7]%[ma:027b84a3]
	Trong chuyến tham quan, một  lớp học muốn thuê một hướng dẫn viên cho chuyến tham quan, có hai công ty đã được liên hệ để lấy thông tin về giá. Công ty $A$ có phí dịch vụ ban đầu là $375\,000$ đồng cộng với $5\,000$ đồng cho mỗi km hướng dẫn. Công ty $B$ có phí dịch vụ ban đầu là $250\,000$ đồng cộng với $7\,500$ đồng cho mỗi km hướng dẫn.
	\choiceTF 
		{\True Lớp học chọn công  ty $B$ sẽ có lợi hơn nếu tổng khoảng cách đi lại là $40$ km}
		{\True Lớp học  chọn công ty $A$ sẽ có lợi hơn nếu tổng khoảng cách đi lại lớn hơn $50$ km}
		{\True Lớp học  chọn công ty $B$ sẽ có lợi hơn nếu tổng khoảng cách đi lại nhỏ hơn $50$ km}
		{Lớp học chọn công ty $B$ sẽ có lợi hơn nếu tổng khoảng cách đi  lại là $60$ km}
	\loigiai{ 
		Gọi $x$ $(x>0)$ là số km lớp đó đi tham quan, khi  đó\\
		Số tiền phải trả cho công ty $A$ là $375\,000+5\,000 x$.\\ 
		Số tiền phải trả cho công ty  $B$ là $250\,000+7\,500 x$.
		\begin{itemchoice}
			\itemch Lớp học chọn công ty $B$ sẽ có lợi  hơn nếu tổng khoảng cách đi lại là $40$ km.\\
			Khi đó  $x=40$ km thì số tiền phải trả cho công ty $A$ là $375\,000+5\,000\cdot 40=575\,000$ đồng, số tiền phải trả cho công ty $B$ là $250\,000+7\,500\cdot 40=550\,000$ đồng.\\
			Vậy chọn công ty $B$ sẽ có  lợi hơn.
			\itemch Lớp học chọn công ty $A$ sẽ có lợi hơn nếu tổng khoảng cách đi lại  lớn hơn $50$ km.\\
			Việc chọn công ty $A$ có lợi hơn nếu số tiền phải trả cho công ty $A$ ít hơn số tiền phải trả cho công ty $B$ tức là $375\,000+5\,000  x<250\,000+7\,500 x \Leftrightarrow x>50$ km.\\
			Vậy  thuê công ty $A$ sẽ có lợi hơn nếu đi với khoảng cách trên $50$ km.
			\itemch Lớp học chọn công ty $B$ sẽ có lợi  hơn nếu tổng khoảng cách đi lại nhỏ hơn $50$ km.\\
			Việc chọn công ty $B$ có lợi hơn nếu số tiền phải trả cho công  ty $B$ ít hơn số tiền phải trả cho công ty $A$ tức là $250\,000+7\,500 x<375\,000+5\,000 x \Leftrightarrow x<50$ km.
			\itemch Lớp học chọn công ty  $B$ sẽ có lợi hơn nếu tổng khoảng cách đi lại là $60$ km.\\
			Khi $x=60$  km thì số tiền phải trả cho công ty $A$ là $375\,000+5\,000\cdot 40=675\,000$ đồng, số tiền phải trả cho công ty $B$ là $250\,000+7\,500\cdot 40=700\,000$ đồng.
		\end{itemchoice}
	} 
\end{ex}

% nguồn: MaHoaNH_DungSai_Lop10
\begin{ex}%[0D3H1-7]%[ma:fd7af878]
	Giả sử giá điện  sinh hoạt trong mỗi tháng dành cho các hộ gia đình được cho bởi bảng sau
	\begin{center}
		\begin{tabular}{|c|c|}
			 \hline
			Mức kWh điện tiêu  thụ & Giá bán điện (VNĐ/ kWh) \\
			 \hline
			Mức 1: từ $0$ đến $100$ kWh &  $1600$\\
			 \hline
			 Mức 2: từ trên $100$ đến $300$ kWh & $2000$\\
			\hline 
			Mức 3: trên  $300$ kWh & $3000$\\
			 \hline
		\end{tabular}
	\end{center}
	Mỗi kết quả dưới đây  đúng hay sai?
	\choiceTF 
		{\True Nếu hộ gia đình tiêu thụ $50$ kWh thì số tiền phải trả là  $80000$ VNĐ}
		{Biểu  thức biểu diễn số tiền phải trả theo số tiêu thụ điện $x$ là $f(x)=2000x+20000$}
		{\True Số tiền điện một  tháng là $420000$ VNĐ thì số điện tiêu thụ là $230$ kWh}
		{Số tiền điện một tháng tối  đa để không vượt quá mức 2 là $500000$ VNĐ}
	 \loigiai{
		\begin{itemchoice}
			\itemch Nếu hộ gia đình tiêu thụ $50$ kWh thì số tiền phải trả là $50\cdot  1600=80000$ VNĐ.
			\itemch 
			Gọi $ x $ kWh là mức điện tiêu thụ. Khi đó, ta có công thức liên hệ giữa số kWh điện tiêu thụ và số tiền tương ứng phải trả $ f(x)  $ như sau
			\allowdisplaybreaks 
			\begin{eqnarray*}
				&& f(x)=\heva{&1600x &(0<x\leq 100)&\\&1600\cdot 100+2000\cdot (x-100) &(100<x\leq 300)&\\&1600\cdot 100+2000\cdot 200+3000\cdot(x-300)&(x>300)&\\}  \\
				&\Leftrightarrow&  f(x)=\heva{&1600x &(0<x\leq 100)&\\&2000x-40000 &(100<x\leq 300)&\\&3000x-340000 &(x>300)&.}
			\end{eqnarray*}
			\itemch 
			Để xác định số tiền điện là $ 420000 $ VNĐ mà hộ gia đình phải trả tương ứng với số kWh điện tiêu thụ bao nhiêu, ta cần xác  định công thức tương ứng.\\
			Với $ f(x)=1600x $,  $ 0<x\leq 100 $ thì $ 0<f(x)\leq 160000 $.\\
			Với $ f(x)=2000x-40000 $, $ 100<x\leq 300 $ thì $ 160000<f(x)\leq  560000 $.\\
			Với $ f(x)=3000x-340000 $,  $ x>300 $ thì $ f(x)>560000 $.\\
			Vì $ 160000<420000<560000 $ nên ứng với số tiền điện $  420000 $ VNĐ thì số kWh điện tiêu thụ của hộ gia đình đó là \[\dfrac{420000+40000}{2000}=230 \text{ kWh.}\]
			Vậy  số kWh điện tiêu thụ của hộ gia đình đó là $ 230 $ kWh.
			\itemch Số tiền điện một tháng tối đa để không vượt quá mức  2 là $560000$ VNĐ.
		\end{itemchoice}
	} 
\end{ex}


%%%%%%%%%%%%%%%% TRẢ LỜI NGẮN %%%%%%%%%%%%%%%%

% nguồn: MaHoaNH_TraLoiNgan_Lop10
\begin{bt}%[0D3H1-7]%[ma:8720ab84]
	Một hãng taxi có bảng giá như  sau (sử dụng cho taxi $4$ chỗ)
	\begin{center}
		\begin{tabular}{|c|c|c|}
			 \hline
			\textbf{Giá mở cửa (0{,}5km)} & \textbf{Giá cước các km tiếp theo} & \textbf{Giá cước từ km thứ  31} \\
			 \hline
			11 000 đồng & 14 500 đồng/km  & 11 600 đồng/km \\
			\hline 
		\end{tabular}
	\end{center}
	Bạn Trân bắt taxi đi quãng đường  $35$ km thì phải trả số tiền là bao nhiêu nghìn đồng (kết quả làm tròn đến hàng nghìn).
	\shortans{497} 
	\loigiai{ 
		Gọi $x$ là số km cần đi $(x\ge 0)$. Gọi $f(x)$ là số tiền phải trả khi  đi $x$ km. Ta có\\
		Nếu $0<x\le 0{,}5$ thì  $f(x)=11\,000$.\\
		Nếu $0{,}5<x\le 30$ thì $f(x)=1\,1000+14\,500(x-0{,}5)=3\,750+14\,500x$.\\ 
		 Nếu $x>30$ thì $f(x)=11\,000+14\,500(30-0{,}5)+11\,600(x-30)=90\,750+11\,600x$.\\
		Vậy ta có công thức hàm  $f(x)$ như sau
		\[f(x)=\heva{&11\,000 &\text{nếu} &&0<x\le  0{,}5 \\&
		3\,750+1\,400x &\text{nếu}  &&0{,}5<x\le 30\\&
		90\,750+11\,600x & \text{nếu} &&x>30.} 
		 \]
		Vì  bạn Trân muốn đi $35$ km nên $x=35>30$ nên\\ $f(35)=90\,750+11\,600\cdot 35=496\,750$ (đồng)\\
		Vậy để đi $35$ km bạn trân phải trả $497$ nghìn đồng. 
	} 
\end{bt}

% nguồn: MaHoaNH_TraLoiNgan_Lop10
\begin{bt}%[0D3V1-7]%[ma:a4df2e41]
	Một cửa hàng nhân dịp Tết đã đồng loạt giảm giá các sản phẩm. Trong đó có chương trình nếu mua một gói kẹo thứ hai trở đi sẽ được giảm $15\%$ so với giá ban đầu. Biết giá gói đầu là $80\,000$ đồng. Bạn An có $600\,000$ đồng. Hỏi bạn An có thể mua tối đa bao nhiêu gói kẹo? 
	 \shortans{8}
	 \loigiai{
		Giả sử bạn An mua được $x$  gói kẹo ($x$ nguyên dương).\\
		Khi đó gói thứ  nhất bạn An phải trả $80\,000$ đồng.\\
		Số gói kẹo còn lại là $x-1$ và bạn An chỉ phải  trả
		\begin{center}
			$80\,000-15\% \cdot 80\,000 =  68\,000$ đồng (mỗi gói).
		\end{center}
		Vậy số tiền phải trả khi mua kẹo được tính theo  công thức
		 $$y=80\,000+(x-1)\cdot 68\,000=68\,000x+12\,000.$$
		Số tiền bạn An dùng mua kẹo phải không quá  $600\,000$ đồng, suy ra
		$$68\,000x+12\,000 \le 600\,000 \Rightarrow x  \le \dfrac{588\,000}{68\,000} \approx 8{,}588.$$
		Vậy với số tiền hiện có, bạn An chỉ  có thể mua được tối đa $8$ gói kẹo.
	 }
\end{bt}


%%%%%%%%%%%%%%%% TỰ LUẬN %%%%%%%%%%%%%%%%

% nguồn: MaHoaNH_TuLuan_Lop10
\begin{bt}%[0D3V1-7]%[ma:1d1f64d2]
	Một công ty dịch vụ cho thuê xe hơi vào dịp tết với giá thuê mỗi chiếc xe hơi như sau: khách thuê tối thiểu phải thuê trọn ba ngày tết (mùng $1$, $2$, $3$) với giá $1\,000\,000$ đồng/ngày; những ngày còn lại (nếu khách còn thuê) sẽ được tính giá thuê là $700\,000$ đồng/ngày. Anh Việt định dành ra một khoản tối đa là $5\,500\,000$ đồng cho phí thuê xe đi chơi trong dịp tết. Hỏi anh Việt có thể thuê xe của công ty trên  tối đa bao nhiêu ngày?
	\loigiai{ 
		 Ta có $T=3\,000\,000+700\,000(x-3)=900\,000+700\,000x$ với $x\ge 3$, $x\in \mathbb{N}$.\\
		Xét bất phương  trình
		 \[900\,000+700\,000x\le 5\,500\,000\Leftrightarrow 9+7x\le 55\Leftrightarrow x\le \dfrac{46}{7} \approx 6{,}57.\]
		Vậy với khoản tiền $5\,500\,000$ đồng, anh Việt chỉ có thể thuê một chiếc xe  của công ty tối đa $6$ ngày.
	 }
\end{bt}

% nguồn: MaHoaNH_TuLuan_Lop10
\begin{bt}%[0D3V1-7]%[ma:1f565bfe]
	Một công ty dịch vụ cho thuê xe hơi vào dịp tết với giá thuê mỗi chiếc xe hơi như sau: khách thuê tối thiểu  phải thuê trọn ba ngày tết (mùng $1$, $2$, $3$) với giá $1\,000\,000$ đồng/ngày; những ngày còn lại (nếu khách còn thuê) sẽ được tính giá thuê là $700\,000$ đồng/ngày. Giả sử $T$ là tổng số tiền mà khách phải trả khi thuê một chiếc xe hơi của công ty và $x$ là số ngày thuê của khách. Khi đó Anh Bình định dành ra một khoản tối đa là $10$ triệu đồng cho phí thuê xe đi chơi trong dịp tết, khi đó anh Bình có thể thuê xe của công ty trên tối đa bao nhiêu ngày?
	\loigiai{ 
		 Ta có $T=3\,000\,000+700\,000(x-3)=900\,000+700\,000x$.\\
		Xét bất phương  trình
		\begin{eqnarray*}
			900\,000+700\,000x\le  10\,000\,000\Leftrightarrow 9+7x\le 100\Leftrightarrow x\le\dfrac{81}{7}\approx 11{,}57.
		\end{eqnarray*}
		Vậy với khoản tiền $10$ triệu đồng, anh  Bình chỉ có thể thuê một chiếc xe tối đa $11$ ngày.
	 }
\end{bt}

% nguồn: MaHoaNH_TuLuan_Lop10
\begin{bt}%[0D3V1-7]%[ma:1d494dcf]
	Một cửa hàng nhân dịp Noel đã đồng loạt giảm giá các sản phẩm. Trong đó có chương trình nếu mua một gói kẹo thứ hai trở đi sẽ được giảm $10\%$ so với giá ban đầu. Biết giá gói đầu là $60\,000$ đồng. Bạn An có $500\,000$ đồng. Hỏi bạn An có thể mua tối đa  bao nhiêu gói kẹo?
	\loigiai{ 
		Giả sử bạn An mua được  $x$ gói kẹo ($x$ nguyên dương).\\
		Khi đó gói  thứ nhất bạn An phải trả $60\,000$ đồng.\\
		Số gói kẹo còn lại là $x-1$ và bạn An chỉ phải  trả
		\begin{center}
			 $60\,000-10\% \cdot 60\,000=54\,000$ đồng (mỗi gói).
		\end{center}
		Vậy số tiền phải trả khi mua kẹo  được tính theo công thức
		$$y=60\,000+(x-1)\cdot 54\,000=54\,000x+6\,000.$$ 
		Số tiền bạn An dùng mua kẹo phải không  quá $500\,000$ đồng, suy ra
		$$54\,000x+6\,000\le 500\,000 \Rightarrow  x\le\dfrac{247}{27}\approx 9{,}148.$$
	Vậy với số  tiền hiện có, bạn An chỉ có thể mua được tối đa $9$ gói kẹo.}
\end{bt}
